\section{Anhang: COLMAP-Vorstudie}

Die COLMAP-Voruntersuchung wurde als separater Benchmark mit COLMAP
\parencite{colmapProject} durchgeführt und ist
nicht mit der späteren GPU-Gesamtpipeline gleichzusetzen. Sie bewertet nur die
Sparse-SfM-Stufe auf dem Alurohr-Testdatensatz.

\begin{table}[H]
\centering
\small
\begin{tabularx}{\textwidth}{Xrrrrr}
\toprule
Konfiguration & Punkte & Punkte/Bild & Tracks & Reprojektion & Zeit (s)\\
\midrule
Plain-SIFT 4096 & 139.449 & 581 & 6,69 & 0,693 Pixel & 876\\
Plain-SIFT 4096 + Guided & 142.208 & 593 & 7,44 & 0,790 Pixel & 1.678\\
Plain-SIFT 8192 & 140.594 & 586 & 6,71 & 0,698 Pixel & 925\\
CRF 18 & 94.616 & 394 & 6,64 & 0,728 Pixel & 677\\
CRF 28 & 93.703 & 390 & 6,43 & 0,761 Pixel & 725\\
\bottomrule
\end{tabularx}
\caption{Optionale COLMAP-SfM-Voruntersuchung. Die Läufe sind ein separater
Benchmark und kein vollständiger STS-/SuGaR-Vergleich.}
\label{tab:colmap-vorstudie}
\end{table}

Die Ergebnisse begründen den Projektarbeitsstandard Plain-SIFT mit 4096
Merkmalen, Sequential-Overlap 15 und deaktiviertem Guided Matching. Guided
Matching brachte nur einen kleinen Punktezuwachs, erhöhte jedoch die Laufzeit
stark und verschlechterte den mittleren Reprojektionsfehler. Die vollständigen
Rohberichte liegen im digitalen Anlagenband unter \path{02_COLMAP_Tests/}
(siehe \texttt{ABGABE\_Index.md}).
